\chapter{Incumbencias del Ingeniero en Sistemas de Información}

La tabla~\ref{tab:incumbencias} relaciona las quince incumbencias del título con lo realizado en el proyecto. Cuando el proyecto no ejercitó una incumbencia, se indica con «No aplica» y el motivo; cuando la ejercitó solo en parte, se aclara el alcance.

{\small
\begin{longtable}{>{\raggedright\arraybackslash}p{0.46\textwidth} >{\raggedright\arraybackslash}p{0.46\textwidth}}
    \caption{Incumbencias aplicadas en el proyecto}\label{tab:incumbencias} \\
    \toprule
    \textbf{Incumbencia} & \textbf{Justificación en el proyecto} \\
    \midrule
    \endfirsthead
    \toprule
    \textbf{Incumbencia} & \textbf{Justificación en el proyecto} \\
    \midrule
    \endhead
    1. Participar en la toma de decisiones estratégicas de una organización y asesorar, en concordancia con las mismas, acerca de las políticas de desarrollo de sistemas de información & Parcial. El equipo definió con el cliente el problema y los objetivos del proyecto (tabla~\ref{tab:objetivos}) y propuso decisiones de política tecnológica, como el procesamiento local sin servicios en la nube (tabla~\ref{tab:decisiones}) y el modelo de implementación y soporte anual (sección~\ref{sec:modelo-negocio}). \\ \midrule
    2. Evaluar, clasificar y seleccionar proyectos de sistemas de información y evaluar y seleccionar alternativas de asistencia interna & Parcial. El equipo comparó las soluciones existentes para la comunicación con personas sordas (sección~\ref{sec:similares}) y las metodologías de desarrollo aplicables (tabla~\ref{tab:metodologias}), y eligió su proyecto entre propuestas de tema (tabla~\ref{tab:hitos}). No clasificó una cartera de proyectos de una organización. \\ \midrule
    3. Planificar, efectuar y evaluar los estudios de factibilidad inherentes a todo proyecto de diseño de sistemas de información y de modificación o reemplazo de los mismos, así como los sistemas de computación asociados & Se estimó la factibilidad económica con el Plan Financiero (tabla~\ref{tab:presupuesto}), la factibilidad temporal con el cronograma (tabla~\ref{tab:hitos}) y los riesgos del proyecto (tabla~\ref{tab:riesgos}). La factibilidad técnica se verificó con el PMV, que cumplió la latencia y la exactitud previstas (tabla~\ref{tab:crit-res}). \\ \midrule
    4. Planificar, dirigir, ejecutar y controlar el relevamiento, análisis, diseño, desarrollo, implementación y prueba de sistemas de información & El proyecto recorrió todo el ciclo: relevamiento y análisis con la ERS y los casos de uso, diseño de arquitectura y de datos, implementación con Docker y pruebas de sistema en dos ciclos (capítulo~\ref{cap:desarrollo}), planificadas y controladas con el plan de gestión (capítulo~\ref{cap:gestion}). \\ \midrule
    5. Evaluar y seleccionar los sistemas de programación disponibles con miras a su utilización en sistemas de información & Se compararon y seleccionaron los lenguajes, bibliotecas y motores del sistema (MediaPipe, PyTorch, FastAPI, Ollama y Piper, entre otros) según los criterios de ejecución local y privacidad (tablas~\ref{tab:tecnologias} y~\ref{tab:comparativa}). \\ \midrule
    6. Evaluar y seleccionar, desde el punto de vista de los sistemas de información, los equipos de procesamiento y comunicación y los sistemas de base & Se definió que la estación opere con una PC, una cámara web y pantallas comerciales, sin hardware especializado (CRIT-06, tabla~\ref{tab:crit}), y se verificó el sistema sobre una PC comercial (tablas~\ref{tab:req-oficiales} y~\ref{tab:especificaciones}). Los riesgos de hardware se trataron en la tabla~\ref{tab:respuesta}. \\ \midrule
    7. Organizar y dirigir el área de sistemas; determinar el perfil de los recursos humanos necesarios y contribuir a su selección y formación & No aplica: el proyecto no incluyó la dirección de un área de sistemas de una organización ni la selección de personal. Solo se organizó el propio equipo, con roles y esfuerzo por paquete de trabajo (tablas~\ref{tab:roles} y~\ref{tab:edt}). \\ \midrule
    8. Participar en la elaboración de programas de capacitación para la utilización de sistemas de información & Se elaboraron el manual de usuario y el de instalación (capítulos~\ref{cap:manual-usuario} y~\ref{cap:manual-instalacion}), y la capacitación del personal se incluyó en el componente de implementación del modelo de ingresos (sección~\ref{sec:modelo-negocio}). Todavía no se dictó capacitación en el sanatorio. \\ \midrule
    9. Determinar y controlar el cumplimiento de las pautas técnicas que rigen el funcionamiento y la utilización de recursos informáticos en cada organización & No aplica en esta etapa: el PMV no se implantó en la infraestructura del cliente, por lo que no se controlaron pautas técnicas de la organización. El manual de instalación (capítulo~\ref{cap:manual-instalacion}) documenta los requisitos de operación del sistema. \\ \midrule
    10. Elaborar métodos y normas a seguir en cuestiones de seguridad y privacidad de la información procesada y la generada por los sistemas de información; participar en la determinación de las acciones a seguir en esta materia y evaluar su aplicación & La privacidad se incorporó desde el diseño: procesamiento local, consentimiento explícito previo a la cámara (CRIT-07) y ausencia de almacenamiento persistente de imágenes (CRIT-04). Su aplicación se evaluó con los casos de prueba de consentimiento y de privacidad y seguridad (tablas~\ref{tab:crit-res} y~\ref{tab:casos}). \\ \midrule
    11. Elaborar métodos y normas a seguir en cuestión de salvaguardia y control de los recursos físicos y lógicos de un sistema de computación; participar en la determinación de las acciones a seguir en esta materia y evaluar su aplicación & Parcial, solo sobre recursos lógicos: se restringieron los orígenes permitidos a la red local, se validaron las entradas y se limitó la información persistida a los reportes y los eventos de auditoría anónimos (tablas~\ref{tab:decisiones} y~\ref{tab:casos}). No se trataron recursos físicos de una organización. \\ \midrule
    12. Desarrollar modelos de simulación, sistemas expertos y otros sistemas informáticos destinados a la resolución de problemas y asesorar en su aplicación & Se desarrolló un sistema de reconocimiento de señas basado en una red neuronal LSTM, con reconstrucción de oraciones y síntesis de voz locales (tabla~\ref{tab:tecnologias}), evaluado en la sección~\ref{sec:res-pruebas}, y se elaboró el manual para su aplicación (capítulo~\ref{cap:manual-usuario}). \\ \midrule
    13. Realizar auditorías en áreas de sistemas y centros de cómputos, así como en los sistemas de información utilizados & No aplica: el proyecto no incluyó auditorías de un área de sistemas ni de sistemas ajenos. Las pruebas de sistema (sección~\ref{sec:modelo-pruebas}) verifican el producto propio, y el registro de eventos anónimos (RF-04) habilita auditorías futuras. \\ \midrule
    14. Realizar arbitrajes, pericias y tasaciones referidas a los sistemas de información y a los medios de procesamiento de datos & No aplica: el proyecto no incluyó arbitrajes, pericias ni tasaciones. El precio de venta estimado (tabla~\ref{tab:presupuesto}) corresponde al propio proyecto y no a una tasación de sistemas de terceros. \\ \midrule
    15. Realizar estudios e investigaciones conducentes a la creación y mejoramiento de técnicas de desarrollo de sistemas de información y nuevas aplicaciones de la tecnología informática existente & Se realizó una investigación del estado del arte (capítulo~\ref{cap:estado-arte}) y se aplicó tecnología existente (visión por computadora y redes de secuencias) a una nueva aplicación: la traducción de un vocabulario clínico de la LSA, con un dataset propio y un modo de captura para ampliarlo. \\
    \bottomrule
\end{longtable}
}
